\documentclass[tikz,border=10pt]{standalone}
\input{../../docs/tikz/dag_preamble.tex}

\begin{document}

% Agent visual validation loop:
% 1. Compile this file to PDF after every substantive DAG edit.
% 2. Inspect the PDF or a rendered PNG, not just the TeX source.
% 3. Increase spacing or reroute arrows until no box, legend, note, edge, or label overlaps.
%
% Intended lifecycle:
% - During paper intake, replace the scaffold below with every named result in the
%   paper and the dependency arrows between them.
% - After intake, treat this as a stable topology: update node styles/status/text
%   as proofs close, but avoid adding boxes or arrows unless the original named
%   result inventory or dependency map was actually incomplete.
%
% Arrow direction rule:
% Draw edges from prerequisite/source node to dependent/target node.  The
% `dag_arrow` and `dag_dashed_arrow` styles use an explicit end-arrow (`->`),
% so the arrowhead should always land on the node that consumes the dependency.
% Do not reverse the coordinate order just to make routing easier; use anchors
% or an orthogonal route through empty space instead.

\begin{tikzpicture}[
  x=1cm,
  y=1cm,
  every node/.append style={outer sep=4pt}
]

\dagPaperMetadata{Simpler Than You Think: The Practical Dynamics of Ranked Choice Voting}{Sanyukta Deshpande; Nikhil Garg; Sheldon H. Jacobson}{Journal of Computational Social Science; arXiv:2602.14329}{2026}{https://link.springer.com/article/10.1007/s42001-026-00470-7}

% --- Legend: README status vocabulary plus formalized node-type styles. ---
\dagPaperLegendRightOfMetadata{
\node[dag_result, dag_template_legend] (legRes) at (0,0) {formalized\\result};
\node[dag_lemma, dag_template_legend] (legLem) at (4.2,0) {formalized\\lemma};
\node[dag_model, dag_template_legend] (legDef) at (8.4,0) {formalized\\definition};
\node[dag_caveat_legend] (legCav) at (12.6,0) {formalized\\with caveat};
\node[dag_partial, dag_template_legend] (legPart) at (0,-2.6) {partially\\formalized};
\node[dag_conditional, dag_template_legend] (legCond) at (4.2,-2.6) {conditional};
\node[dag_scaffold, dag_template_legend] (legScaf) at (8.4,-2.6) {scaffold};
\node[dag_unformalized, dag_template_legend] (legNot) at (12.6,-2.6) {not started /\\not formalized};
}
\daglegend{(legRes)(legLem)(legDef)(legCav)(legPart)(legCond)(legScaf)(legNot)}{Legend}

\begin{dagPaperBody}

\node[dag_model] (BallotLayer) at (0,-3.8) {
  \textbf{Ballot and support layer} \\
  Prefix/suffix robustness, exhausted prefixes, and active-support counts
};
\node[dag_result] (Prop1) at (6.4,-3.8) {
  \textbf{Proposition 1} \\
  Robust strategy transforms over DGJ24 support-count Algorithm A outputs
};
\node[dag_result] (Prop2) at (12.8,-3.8) {
  \textbf{Proposition 2} \\
  Exhausted-ballot viable-set formulas and count-test closeout
};

\node[dag_model] (RemovalLayer) at (0,-8.4) {
  \textbf{Candidate-removal layer} \\
  Strict support, Algorithm 2 condition, and Algorithm 3 extensions
};
\node[dag_result] (Thm21) at (6.4,-8.4) {
  \textbf{Theorem 2.1} \\
  Full Algorithm 3 source branch and one-survival step-facts route
};

\node[dag_result] (Thm22) at (12.8,-11.0) {
  \textbf{Theorem 2.2} \\
  Algorithm 4 pairwise condition and containment output theorem
};

\draw[dag_arrow] (BallotLayer.east) -- (Prop1.west);
\draw[dag_arrow] (BallotLayer.east) -- (Prop2.west);
\draw[dag_arrow] (RemovalLayer.east) -- (Thm21.west);
\draw[dag_arrow] (RemovalLayer.south east) |- ($(Thm22.west)+(-0.5,0)$) -- (Thm22.west);
\end{dagPaperBody}
\end{tikzpicture}
\end{document}
